% 사물인터넷과 빅데이터 - 2주차 강의자료
% 바이브코딩 소개 & 센서 탐색
% XeLaTeX로 컴파일: xelatex 강의자료_2주차.tex

\documentclass[11pt, a4paper]{article}

% 한국어 지원
\usepackage{fontspec}
\setmainfont{Noto Sans CJK KR}[
  UprightFont = *-Regular,
  BoldFont = *-Bold
]

% 페이지 설정
\usepackage[top=2.5cm, bottom=2.5cm, left=2.5cm, right=2.5cm]{geometry}
\usepackage{setspace}
\onehalfspacing

% 색상 및 박스
\usepackage[dvipsnames]{xcolor}
\usepackage{tcolorbox}
\tcbuselibrary{skins, breakable}

% 표 및 그래픽
\usepackage{booktabs}
\usepackage{tabularx}
\usepackage{graphicx}
\usepackage{float}
\usepackage{wrapfig}

% TikZ 그래픽
\usepackage{tikz}
\usetikzlibrary{shapes.geometric, arrows.meta, positioning, calc, fit, backgrounds, shadows}

% 목록 및 링크
\usepackage{enumitem}
\usepackage{hyperref}
\hypersetup{
    colorlinks=true,
    linkcolor=NavyBlue,
    urlcolor=RoyalBlue,
    bookmarksnumbered=true
}

% 헤더/푸터
\usepackage{fancyhdr}
\pagestyle{fancy}
\fancyhf{}
\fancyhead[L]{\small\textcolor{gray}{사물인터넷과 빅데이터}}
\fancyhead[R]{\small\textcolor{gray}{2주차: 바이브코딩 소개 \& 센서 탐색}}
\fancyfoot[C]{\thepage}
\renewcommand{\headrulewidth}{0.4pt}

% 섹션 스타일
\usepackage{titlesec}
\titleformat{\section}
    {\Large\bfseries\color{NavyBlue}}
    {\thesection}{1em}{}[\titlerule]
\titleformat{\subsection}
    {\large\bfseries\color{RoyalBlue}}
    {\thesubsection}{1em}{}

% 커스텀 박스 정의
\newtcolorbox{keypoint}[1][]{
  enhanced,
  colback=blue!5,
  colframe=NavyBlue,
  fonttitle=\bfseries,
  title={💡 핵심 포인트},
  #1
}

\newtcolorbox{example}[1][]{
  enhanced,
  colback=green!5,
  colframe=ForestGreen,
  fonttitle=\bfseries,
  title={📝 예시},
  #1
}

\newtcolorbox{warning}[1][]{
  enhanced,
  colback=orange!5,
  colframe=BurntOrange,
  fonttitle=\bfseries,
  title={⚠️ 주의},
  #1
}

\newtcolorbox{promptbox}[1][]{
  enhanced,
  colback=gray!5,
  colframe=gray!50,
  fonttitle=\bfseries,
  title={💬 프롬프트 예시},
  #1
}

% 색상 정의
\definecolor{mainblue}{RGB}{0, 102, 204}
\definecolor{accentgreen}{RGB}{40, 167, 69}
\definecolor{lightgray}{RGB}{245, 245, 245}
\definecolor{warningorange}{RGB}{255, 152, 0}
\definecolor{microbitcolor}{RGB}{0, 150, 100}

% 문서 정보
\title{
  \vspace{-2cm}
  {\Huge\bfseries\color{NavyBlue} 사물인터넷과 빅데이터}\\[0.5em]
  {\LARGE 2주차: 바이브코딩 소개 \& 센서 탐색}
}
\author{
  양창모\\
  청주교육대학교 컴퓨터교육과
}
\date{2026학년도 1학기}

\begin{document}

% 표지
\maketitle
\thispagestyle{empty}

\vfill

\begin{center}
\begin{tcolorbox}[width=0.8\textwidth, colback=blue!5, colframe=NavyBlue]
\centering
\large
\textbf{학습 목표}

\vspace{0.5em}

\begin{itemize}[leftmargin=2em]
  \item 바이브코딩의 개념을 이해한다
  \item Claude를 활용하여 간단한 웹페이지를 만든다
  \item 마이크로비트 내장 센서와 Grove 확장 센서를 탐색한다
  \item 프로젝트 아이디어를 구상한다
\end{itemize}
\end{tcolorbox}
\end{center}

\vfill

\begin{center}
\small
\textit{청주교육대학교 컴퓨터교육과}\\
\textit{v1.1.0.0 | 2026-01-28}
\end{center}

\newpage

% 목차
\tableofcontents
\newpage

%===========================================
\section{지난 주 복습}
%===========================================

\subsection{1주차 과제 확인}

지난 주에는 마이크로비트를 처음 접하고 LED에 자기 이름을 표시하는 과제를 수행했습니다.

\begin{keypoint}
\textbf{점검 사항}
\begin{itemize}
  \item MakeCode 에디터 접속 성공 여부
  \item 마이크로비트 연결 및 프로그래밍 성공 여부
  \item LED 문자 표시 원리 이해 여부
\end{itemize}
\end{keypoint}

\subsection{마이크로비트란?}

마이크로비트는 영국 BBC에서 교육용으로 개발한 초소형 컴퓨터입니다. 현재 V2 버전이 출시되어 더 많은 기능을 제공합니다.

\begin{figure}[H]
\centering
\includegraphics[width=0.6\textwidth]{microbitdrawing2bluemicrophoneon.png}
\caption{마이크로비트 V2 앞면 - LED, 버튼, 마이크 등}
\end{figure}

\subsubsection{주요 구성 요소}

\textbf{앞면 (Front)}
\begin{itemize}
  \item \textbf{5×5 LED 디스플레이}: 25개 LED로 문자, 숫자, 그림 표시
  \item \textbf{버튼 A, B}: 입력 장치
  \item \textbf{터치 로고}: V2에서 추가된 터치 센서
  \item \textbf{마이크}: V2에서 추가된 소리 감지 센서
  \item \textbf{엣지 커넥터}: 외부 장치 연결용 핀 (0, 1, 2, 3V, GND 등)
\end{itemize}

\begin{figure}[H]
\centering
\includegraphics[width=0.6\textwidth]{microbitdrawingback13.png}
\caption{마이크로비트 V2 뒷면 - 프로세서, 센서, 연결부}
\end{figure}

\textbf{뒷면 (Back)}
\begin{itemize}
  \item \textbf{프로세서}: 프로그램을 실행하는 CPU
  \item \textbf{가속도 센서}: 움직임, 기울기 감지
  \item \textbf{나침반 (자기 센서)}: 방위 측정
  \item \textbf{USB 포트}: 프로그래밍 및 전원 공급
  \item \textbf{배터리 커넥터}: 휴대용 전원 연결
  \item \textbf{리셋 버튼}: 프로그램 재시작
\end{itemize}

\begin{figure}[H]
\centering

\includegraphics[width=0.15\textwidth]{v2only.png}
\caption{V2 전용 기능 표시}
\end{figure}

\subsection{MakeCode 에디터 복습}

MakeCode는 마이크로비트를 프로그래밍할 수 있는 블록 코딩 환경입니다.

\textbf{주요 영역}
\begin{enumerate}
  \item \textbf{시뮬레이터}: 코드 실행 결과를 바로 확인
  \item \textbf{블록 팔레트}: 사용 가능한 블록들 (기본, 입력, LED, 루프 등)
  \item \textbf{코딩 영역}: 블록을 끌어다 놓아 프로그램 작성
\end{enumerate}

\begin{keypoint}
\textbf{기본 블록 구조}

모든 마이크로비트 프로그램은 다음 블록으로 시작합니다:
\begin{itemize}
  \item \textbf{시작하면}: 전원이 켜지거나 리셋할 때 한 번 실행
  \item \textbf{무한반복}: 계속 반복해서 실행
\end{itemize}
\end{keypoint}

\newpage

%===========================================
\section{바이브코딩이란?}
%===========================================

\subsection{바이브코딩의 개념}

\begin{tcolorbox}[enhanced, colback=blue!10, colframe=mainblue, 
                  arc=5mm, fontupper=\Large\bfseries, center title]
"코드를 짜지 않고 코드를 만든다"
\end{tcolorbox}

\vspace{1em}

바이브코딩(Vibe Coding)은 \textbf{AI와 대화하며 프로그램을 만드는} 새로운 프로그래밍 방식입니다. 전통적인 코딩에서는 문법을 배우고 직접 코드를 작성해야 했지만, 바이브코딩에서는 \textbf{자연어로 원하는 것을 설명}하면 AI가 코드를 생성해줍니다.

\subsection{전통적 프로그래밍 vs 바이브코딩}

\begin{table}[H]
\centering
\begin{tabularx}{\textwidth}{|>{\centering\arraybackslash}X|>{\centering\arraybackslash}X|}
\hline
\rowcolor{lightgray}
\textbf{전통적 프로그래밍} & \textbf{바이브코딩} \\
\hline
문법 학습 필수 & 문법 학습 불필요 \\
\hline
직접 코드 작성 & AI와 대화로 작성 \\
\hline
오류 수정에 시간 소요 & AI가 즉시 수정 \\
\hline
시간 소요 +++ & 시간 소요 + \\
\hline
난이도 높음 & 난이도 낮음 \\
\hline
\textbf{코드 작성 능력 중요} & \textbf{문제 정의 능력 중요} \\
\hline
\end{tabularx}
\caption{전통적 프로그래밍과 바이브코딩 비교}
\end{table}

\subsubsection{전통적 프로그래밍 프로세스}

\begin{center}
\begin{tikzpicture}[
  node distance=1cm,
  box/.style={rectangle, rounded corners=3pt, minimum width=2cm, minimum height=0.8cm, 
              text centered, draw=black, fill=lightgray, font=\small},
  arrow/.style={-Stealth, thick, draw=mainblue}
]

\node[box] (p1) {사람};
\node[box, right=of p1] (p2) {문법 학습};
\node[box, right=of p2] (p3) {코드 작성};
\node[box, right=of p3] (p4) {디버깅};
\node[box, right=of p4, fill=accentgreen!20] (p5) {프로그램};

\draw[arrow] (p1) -- (p2);
\draw[arrow] (p2) -- (p3);
\draw[arrow] (p3) -- (p4);
\draw[arrow] (p4) -- (p5);

% 오류 순환
\draw[-Stealth, thick, draw=warningorange, dashed, bend left=45] 
  (p4.north) to node[above, font=\tiny] {오류 발생} (p2.north);

\end{tikzpicture}
\end{center}

\subsubsection{바이브코딩 프로세스}

\begin{center}
\begin{tikzpicture}[
  node distance=1cm,
  box/.style={rectangle, rounded corners=3pt, minimum width=2cm, minimum height=0.8cm, 
              text centered, draw=black, fill=lightgray, font=\small},
  highlight/.style={rectangle, rounded corners=3pt, minimum width=2cm, minimum height=0.8cm,
                    text centered, draw=accentgreen, fill=accentgreen!20, line width=1.5pt, font=\small},
  arrow/.style={-Stealth, thick, draw=mainblue}
]

\node[box] (v1) {사람};
\node[box, right=of v1] (v2) {문제 정의};
\node[highlight, right=of v2] (v3) {AI 대화};
\node[box, right=of v3] (v4) {결과 검증};
\node[highlight, right=of v4] (v5) {프로그램};

\draw[arrow] (v1) -- (v2);
\draw[arrow] (v2) -- (v3);
\draw[arrow] (v3) -- (v4);
\draw[arrow] (v4) -- (v5);

% 수정 순환
\draw[-Stealth, thick, draw=accentgreen, dashed, bend right=45] 
  (v4.south) to node[below, font=\tiny] {수정 요청} (v2.south);

\end{tikzpicture}
\end{center}

\newpage

\subsection{왜 바이브코딩인가?}

\subsubsection{초등교사에게 필요한 역량}

초등교사는 학생들에게 \textbf{프로그래밍 문법}을 가르치는 것보다 \textbf{컴퓨팅 사고력}과 \textbf{문제 해결 능력}을 길러주는 것이 더 중요합니다.

\begin{keypoint}
\textbf{바이브코딩으로 기를 수 있는 역량}
\begin{enumerate}
  \item \textbf{문제 정의 능력}: "무엇을 만들고 싶은가?"를 명확히 하기
  \item \textbf{비판적 사고}: "이 결과가 맞나? 확인해보자"
  \item \textbf{AI 협업 능력}: AI에게 효과적으로 요청하는 방법
  \item \textbf{컴퓨팅 사고력}: 문제를 단계별로 나누어 해결하기
\end{enumerate}
\end{keypoint}

\subsubsection{2022 개정 교육과정과의 연계}

2022 개정 교육과정에서는 초등학생의 \textbf{SW/AI 교육}이 강조됩니다. 바이브코딩은 다음과 같이 교육과정과 연계됩니다:

\begin{table}[H]
\centering
\begin{tabularx}{\textwidth}{|l|X|}
\hline
\rowcolor{lightgray}
\textbf{교육과정 요소} & \textbf{바이브코딩 연계} \\
\hline
실과 - 피지컬 컴퓨팅 & 마이크로비트 + 센서 활용 \\
\hline
실과 - 데이터 수집·분석 & 센서 데이터 수집 및 시각화 \\
\hline
SW/AI 교육 & AI 협업, 문제 해결 \\
\hline
컴퓨팅 사고력 & 문제 정의, 추상화, 분해 \\
\hline
\end{tabularx}
\end{table}

\subsection{바이브코딩의 핵심}

\begin{center}
\begin{tikzpicture}[
  node distance=2cm,
  mainnode/.style={circle, minimum size=2.5cm, text centered, text width=2cm,
                   draw=accentgreen, fill=accentgreen!10, line width=2pt, font=\large\bfseries},
  subnode/.style={circle, minimum size=2cm, text centered, text width=1.8cm,
                  draw=mainblue, fill=mainblue!10, line width=1.5pt, font=\small\bfseries}
]

% 중앙
\node[mainnode] (center) at (0,0) {AI와\\협업};

% 주변 노드들 (시계 방향)
\node[subnode] (define) at (0,2.8) {문제\\정의};
\node[subnode] (converse) at (2.4,1.4) {대화\\능력};
\node[subnode] (verify) at (2.4,-1.4) {결과\\검증};
\node[subnode] (refine) at (0,-2.8) {반복\\수정};
\node[subnode] (iterate) at (-2.4,-1.4) {비판적\\사고};
\node[subnode] (complete) at (-2.4,1.4) {완성};

% 화살표
\foreach \from/\to in {define/converse, converse/verify, verify/refine, 
                       refine/iterate, iterate/complete, complete/define} {
  \draw[-Stealth, line width=1.5pt, draw=accentgreen] (\from) -- (\to);
}

% 중앙 연결
\foreach \node in {define, converse, verify, refine, iterate, complete} {
  \draw[-Stealth, line width=1pt, draw=mainblue, dashed] (\node) -- (center);
}

\end{tikzpicture}
\end{center}

\begin{warning}
\textbf{바이브코딩의 오해}

바이브코딩이 "모든 것을 AI가 해준다"는 의미는 \textbf{아닙니다}. 가장 중요한 것은:
\begin{itemize}
  \item 문제를 \textbf{명확히 정의}하는 것
  \item 결과를 \textbf{비판적으로 검증}하는 것
  \item 만족할 때까지 \textbf{반복 수정}하는 것
\end{itemize}
\end{warning}

\newpage

%===========================================
\section{Claude 사용법}
%===========================================

\subsection{Claude란?}

\textbf{Claude}는 Anthropic에서 개발한 AI 어시스턴트로, 대화를 통해 코드를 생성하고 문제를 해결할 수 있습니다.

\subsubsection{Claude의 특징}

\begin{itemize}
  \item \textbf{긴 대화 컨텍스트 지원}: 이전 대화 내용을 기억
  \item \textbf{코드 생성 및 분석 특화}: HTML, JavaScript, Python 등
  \item \textbf{단계적 수정 가능}: 결과를 보고 수정 요청
  \item \textbf{교육용으로 적합}: 친절한 설명과 예시 제공
\end{itemize}

\textbf{접속 방법}
\begin{itemize}
  \item 웹사이트: \url{https://claude.ai}
  \item 계정 생성 (무료)
  \item 바로 대화 시작
\end{itemize}

\begin{keypoint}
\textbf{다른 AI도 사용 가능합니다}

ChatGPT, Copilot, Gemini 등 다른 AI도 바이브코딩에 사용할 수 있습니다. 이 강의에서는 Claude를 예시로 사용합니다.
\end{keypoint}

\subsection{효과적인 프롬프트 작성법}

\textbf{프롬프트(Prompt)}는 AI에게 보내는 요청이나 질문입니다. 프롬프트를 잘 작성하면 원하는 결과를 더 빠르게 얻을 수 있습니다.

\subsubsection{프롬프트 작성 원칙}

\begin{enumerate}
  \item \textbf{구체적으로 요청하기}
  \begin{itemize}
    \item ❌ 나쁜 예: "웹페이지 만들어줘"
    \item ✅ 좋은 예: "자기소개 웹페이지를 만들어줘. 이름, 전공, 취미를 포함하고 파란색 테마로"
  \end{itemize}
  
  \item \textbf{형식 지정하기}
  \begin{itemize}
    \item "HTML 파일 하나로 만들어줘"
    \item "CSS는 파일 안에 포함해줘"
    \item "설명은 주석으로 넣어줘"
  \end{itemize}
  
  \item \textbf{단계적으로 수정하기}
  \begin{itemize}
    \item 일단 기본 버전을 만들고
    \item 결과를 확인한 후
    \item 수정할 부분을 구체적으로 요청
  \end{itemize}
  
  \item \textbf{예시 제공하기}
  \begin{itemize}
    \item 원하는 스타일의 예시 제시
    \item 데이터 형식 예시 제공
    \item 비슷한 결과물 참고 자료
  \end{itemize}
\end{enumerate}

\newpage

\subsection{프롬프트 예시}

\begin{promptbox}
\textbf{예시 1: 자기소개 웹페이지}

\begin{verbatim}
자기소개 웹페이지를 만들어줘.

이름: 김철수
전공: 수학교육과
취미: 등산, 독서

심플한 디자인으로 HTML 파일 하나로 만들어줘.
배경색은 연한 파란색으로 해줘.
\end{verbatim}
\end{promptbox}

\vspace{1em}

\begin{promptbox}
\textbf{예시 2: 수정 요청}

\begin{verbatim}
방금 만든 웹페이지에서:

1. 배경색을 연한 파란색으로 바꿔줘
2. 글자 크기를 조금 키워줘
3. 사진을 넣을 수 있는 공간을 추가해줘
\end{verbatim}
\end{promptbox}

\vspace{1em}

\begin{promptbox}
\textbf{예시 3: 데이터 시각화}

\begin{verbatim}
다음 데이터를 그래프로 시각화하는 웹페이지를 만들어줘.

시간, 온도, 소음
09:00, 22, 45
10:00, 24, 78
11:00, 26, 120

선 그래프로 보여주고, HTML 한 파일로 만들어줘.
\end{verbatim}
\end{promptbox}

\subsection{프롬프트 작성 팁}

\begin{keypoint}
\textbf{효과적인 프롬프트의 특징}
\begin{itemize}
  \item \textbf{명확함}: 무엇을 원하는지 분명히 표현
  \item \textbf{구체성}: 세부 사항까지 포함
  \item \textbf{맥락 제공}: 어떤 목적인지 설명
  \item \textbf{제약 조건}: 파일 형식, 길이, 스타일 등
  \item \textbf{예시 포함}: 원하는 결과의 예시
\end{itemize}
\end{keypoint}

\begin{warning}
\textbf{주의사항}
\begin{itemize}
  \item 한 번에 너무 많은 것을 요청하지 말기
  \item 막연한 표현 피하기 ("예쁘게", "멋있게" 등)
  \item 불가능한 요청 피하기
  \item 대화 기록을 저장해두기 (과제 제출용)
\end{itemize}
\end{warning}

\newpage

%===========================================
\section{바이브코딩 실습}
%===========================================

\subsection{실습 목표}

Claude와 대화하여 \textbf{자신만의 자기소개 웹페이지}를 만들어봅니다.

\subsection{실습 준비}

\subsubsection{필요한 것}
\begin{itemize}
  \item 컴퓨터 (인터넷 연결)
  \item 웹 브라우저 (Chrome, Edge, Safari 등)
  \item Claude 계정 (\url{https://claude.ai})
  \item 텍스트 에디터 (메모장, VS Code 등)
\end{itemize}

\subsubsection{실습 순서}

\begin{enumerate}
  \item Claude 접속 (\url{https://claude.ai})
  \item 기본 웹페이지 요청 (프롬프트 작성)
  \item Claude가 생성한 코드 복사
  \item 텍스트 에디터에 붙여넣기
  \item 파일 저장 (*.html)
  \item 브라우저에서 열어보기
  \item 수정 요청 (색상, 레이아웃, 내용 등)
  \item 최종 결과물 저장
\end{enumerate}

\subsection{실습 체크리스트}

\begin{table}[H]
\centering
\begin{tabularx}{\textwidth}{|l|X|}
\hline
\rowcolor{lightgray}
\textbf{항목} & \textbf{내용} \\
\hline
\multicolumn{2}{|l|}{\textbf{필수 포함 내용}} \\
\hline
☑ 이름 & 본인의 이름 표시 \\
\hline
☑ 전공 & 전공 또는 학년 정보 \\
\hline
☑ 관심사 & 취미, 관심 분야 등 \\
\hline
\multicolumn{2}{|l|}{\textbf{선택 추가 내용}} \\
\hline
☐ 사진 & 프로필 사진 또는 이미지 \\
\hline
☐ 연락처 & 이메일, SNS 등 \\
\hline
☐ 명언 & 좋아하는 명언이나 좌우명 \\
\hline
☐ 기타 & 자유롭게 추가 \\
\hline
\end{tabularx}
\end{table}

\begin{warning}
\textbf{실습 시 주의사항}
\begin{itemize}
  \item \textbf{대화 기록 저장하기}: 과제 제출에 필요
  \item \textbf{최종 코드 백업하기}: 여러 곳에 저장
  \item \textbf{막히면 질문하기}: 교수님 또는 친구에게
\end{itemize}
\end{warning}

\subsection{실습 예시}

\begin{example}
\textbf{실습 과정 예시}

\textbf{1단계: 첫 프롬프트}
\begin{verbatim}
자기소개 웹페이지를 만들어줘.
이름: 홍길동, 전공: 컴퓨터교육, 취미: 코딩
\end{verbatim}

\textbf{2단계: 결과 확인 후 수정}
\begin{verbatim}
배경색을 하늘색으로 바꾸고,
제목을 크게 만들어줘
\end{verbatim}

\textbf{3단계: 추가 요청}
\begin{verbatim}
사진을 넣을 수 있는 원형 프레임을 추가해줘
\end{verbatim}
\end{example}

\newpage

%===========================================
\section{센서 탐색}
%===========================================

\subsection{마이크로비트 센서 시스템}

마이크로비트는 \textbf{내장 센서}와 \textbf{Grove 확장 센서}를 사용하여 다양한 데이터를 수집할 수 있습니다.

\subsubsection{센서로 무엇을 할 수 있나?}

\begin{itemize}
  \item \textbf{환경 측정}: 온도, 습도, 기압, 공기질 등
  \item \textbf{상태 감지}: 빛, 소리, 움직임 등
  \item \textbf{거리 측정}: 사물이나 사람까지의 거리
  \item \textbf{색상 인식}: 물체의 색상 파악
  \item \textbf{데이터 기록}: 측정한 데이터를 SD카드에 저장
\end{itemize}

\subsection{내장 센서 (6개)}

마이크로비트 V2에는 \textbf{6개의 센서가 기본 내장}되어 있습니다.

\begin{table}[H]
\centering
\begin{tabularx}{\textwidth}{|l|l|X|}
\hline
\rowcolor{lightgray}
\textbf{센서} & \textbf{측정 데이터} & \textbf{특징} \\
\hline
가속도계 & X축, Y축, Z축 & 기울기, 흔들림, 낙하 감지 \\
\hline
나침반 & 방위각 (0°~360°) & 방향 측정 \\
\hline
온도센서 & 온도 (대략적) & CPU 기반, 실제보다 5~10°C 높음 \\
\hline
빛센서 & 밝기 (0~255) & LED를 역이용한 센서 \\
\hline
마이크 (V2) & 소리 크기 (0~255) & V2에서 추가 \\
\hline
터치 로고 (V2) & 터치 이벤트 & V2에서 추가, 정전식 센서 \\
\hline
\end{tabularx}
\caption{마이크로비트 V2 내장 센서}
\end{table}

\begin{keypoint}
\textbf{내장 센서의 장점}
\begin{itemize}
  \item \textbf{별도 연결 불필요}: 마이크로비트만으로 사용 가능
  \item \textbf{프로그래밍 간단}: MakeCode에서 쉽게 사용
  \item \textbf{비용 절감}: 추가 비용 없음
  \item \textbf{교육용 적합}: 기초 센서 학습에 충분
\end{itemize}
\end{keypoint}

\subsection{Grove 확장 센서 (10개)}

더 다양하고 정밀한 측정을 위해 \textbf{Grove 센서}를 추가로 연결할 수 있습니다.

\begin{figure}[H]
\centering
\includegraphics[width=0.4\textwidth]{grovetemperaturehumiditysensordht11preview.webp}
\caption{Grove 온습도 센서 예시}
\end{figure}

\begin{table}[H]
\centering
\begin{tabularx}{\textwidth}{|l|c|l|X|}
\hline
\rowcolor{lightgray}
\textbf{센서} & \textbf{I2C 주소} & \textbf{측정 데이터} & \textbf{활용 예시} \\
\hline
SHT40 온습도 & 0x44 & 정밀온도, 습도 & 쾌적한 교실 환경 측정 \\
\hline
BMP280 기압 & 0x76 & 기압, 고도 & 날씨 변화 예측 \\
\hline
OLED 디스플레이 & 0x3C & 화면출력 & 측정값 실시간 표시 \\
\hline
LIS3DHTR 가속도 & 0x18 & 정밀 X, Y, Z축 & 세밀한 움직임 감지 \\
\hline
TCS34725 컬러 & 0x29 & R, G, B, 밝기 & 물체 색상 인식 \\
\hline
TSL2561 조도 & 0x39 & 럭스 (lux) & 정밀한 조도 측정 \\
\hline
SGP30 공기질 & 0x58 & CO₂, TVOC & 교실 환기 상태 확인 \\
\hline
VL53L0X 거리 & 0x29 & 거리 (mm) & 사람 감지, 줄 서기 측정 \\
\hline
MPR121 터치 & 0x5A & 12채널 터치 & 다양한 터치 입력 \\
\hline
PAJ7620 제스처 & 0x73 & 제스처 패턴 & 손동작 인식 \\
\hline
\end{tabularx}
\caption{Grove 확장 센서 목록}
\end{table}

\newpage

\subsubsection{Grove 센서 상세 소개}

각 센서의 외형과 특징을 자세히 알아봅니다.

\paragraph{1. TSL2561 조도 센서}

\begin{figure}[H]
\centering
\includegraphics[width=0.3\textwidth]{sensor_01_light.jpg}
\caption{TSL2561 조도 센서}
\end{figure}

\begin{itemize}
  \item \textbf{I2C 주소}: 0x39
  \item \textbf{측정 데이터}: 조도 (0.1 ~ 40,000 lux)
  \item \textbf{특징}: 사람의 눈과 유사한 빛 감지, 적외선 보정 기능
  \item \textbf{활용}: 교실 조명 상태 측정, 창가/복도 밝기 비교
\end{itemize}

\paragraph{2. TCS34725 컬러 센서}

\begin{figure}[H]
\centering
\includegraphics[width=0.3\textwidth]{sensor_02_color.jpg}
\caption{TCS34725 컬러 센서}
\end{figure}

\begin{itemize}
  \item \textbf{I2C 주소}: 0x29
  \item \textbf{측정 데이터}: RGB 색상값 (0~255), 밝기
  \item \textbf{특징}: 적색, 녹색, 청색 필터 내장, 백색 LED 포함
  \item \textbf{활용}: 물체 색상 인식, 색깔 분류 게임, 미술 작품 분석
\end{itemize}

\paragraph{3. SHT35/SHT40 온습도 센서}

\begin{figure}[H]
\centering
\includegraphics[width=0.3\textwidth]{sensor_03_sht35.jpg}
\caption{SHT35/SHT40 온습도 센서}
\end{figure}

\begin{itemize}
  \item \textbf{I2C 주소}: 0x44
  \item \textbf{측정 데이터}: 정밀 온도 (-40~125°C), 습도 (0~100\%)
  \item \textbf{특징}: 높은 정확도 (±0.2°C, ±2\%RH), 빠른 응답 속도
  \item \textbf{활용}: 교실 환경 모니터링, 쾌적도 분석, 환기 시점 파악
\end{itemize}

\paragraph{4. LIS3DHTR 가속도 센서}

\begin{figure}[H]
\centering
\includegraphics[width=0.3\textwidth]{sensor_04_accel_gyro.jpg}
\caption{LIS3DHTR 3축 가속도 센서}
\end{figure}

\begin{itemize}
  \item \textbf{I2C 주소}: 0x18
  \item \textbf{측정 데이터}: 3축 가속도 (X, Y, Z)
  \item \textbf{특징}: 내장 센서보다 정밀, 충격 감지 가능
  \item \textbf{활용}: 체육 활동량 측정, 낙상 감지, 진동 분석
\end{itemize}

\paragraph{5. VL53L0X 거리 센서 (ToF)}

\begin{figure}[H]
\centering
\includegraphics[width=0.3\textwidth]{sensor_05_tof.jpg}
\caption{VL53L0X ToF 거리 센서}
\end{figure}

\begin{itemize}
  \item \textbf{I2C 주소}: 0x29
  \item \textbf{측정 데이터}: 거리 (30~2000 mm)
  \item \textbf{특징}: 레이저 기반 정밀 측정, 빛 조건 영향 적음
  \item \textbf{활용}: 사람 감지, 급식실 줄 길이 측정, 장애물 회피
\end{itemize}

\paragraph{6. BMP280 기압 센서}

\begin{figure}[H]
\centering
\includegraphics[width=0.3\textwidth]{sensor_06_bmp280.jpg}
\caption{BMP280 기압 센서}
\end{figure}

\begin{itemize}
  \item \textbf{I2C 주소}: 0x76
  \item \textbf{측정 데이터}: 기압 (300~1100 hPa), 온도
  \item \textbf{특징}: 고도 계산 가능, 날씨 변화 감지
  \item \textbf{활용}: 날씨 예측, 층간 고도차 측정, 기압 변화 관찰
\end{itemize}

\paragraph{7. BME680 복합 환경 센서}

\begin{figure}[H]
\centering
\includegraphics[width=0.3\textwidth]{sensor_07_bme680.jpg}
\caption{BME680 복합 환경 센서}
\end{figure}

\begin{itemize}
  \item \textbf{I2C 주소}: 0x76
  \item \textbf{측정 데이터}: 온도, 습도, 기압, 공기질 (IAQ)
  \item \textbf{특징}: 4가지 센서 통합, VOC 가스 감지
  \item \textbf{활용}: 종합 환경 모니터링, 공기질 지수 계산
\end{itemize}

\paragraph{8. OLED 디스플레이}

\begin{figure}[H]
\centering
\includegraphics[width=0.3\textwidth]{sensor_08_oled.jpg}
\caption{OLED 디스플레이 (128x64)}
\end{figure}

\begin{itemize}
  \item \textbf{I2C 주소}: 0x3C
  \item \textbf{기능}: 텍스트/그래픽 출력 (128×64 픽셀)
  \item \textbf{특징}: 선명한 화면, 넓은 시야각, 저전력
  \item \textbf{활용}: 측정값 실시간 표시, 메시지 전달, 그래프 표시
\end{itemize}

\paragraph{9. MCP9808 정밀 온도 센서}

\begin{figure}[H]
\centering
\includegraphics[width=0.3\textwidth]{sensor_09_mcp9808.jpg}
\caption{MCP9808 정밀 온도 센서}
\end{figure}

\begin{itemize}
  \item \textbf{I2C 주소}: 0x18
  \item \textbf{측정 데이터}: 정밀 온도 (-40~125°C)
  \item \textbf{특징}: 매우 높은 정확도 (±0.25°C), 알람 기능
  \item \textbf{활용}: 정밀 온도 실험, 온도 변화 감지, 냉난방 효율 분석
\end{itemize}

\paragraph{10. AHT20 온습도 센서}

\begin{figure}[H]
\centering
\includegraphics[width=0.3\textwidth]{sensor_10_aht20.png}
\caption{AHT20 온습도 센서}
\end{figure}

\begin{itemize}
  \item \textbf{I2C 주소}: 0x38
  \item \textbf{측정 데이터}: 온도 (-40~85°C), 습도 (0~100\%)
  \item \textbf{특징}: 소형 크기, 낮은 전력 소비, 빠른 응답
  \item \textbf{활용}: 휴대용 환경 측정기, 장기간 데이터 수집
\end{itemize}

\begin{keypoint}
\textbf{센서 선택 팁}
\begin{itemize}
  \item \textbf{내장 센서로 충분한 경우}: 간단한 측정, 교육용 데모
  \item \textbf{정밀도가 필요한 경우}: Grove 확장 센서 사용
  \item \textbf{여러 데이터가 필요한 경우}: BME680 같은 복합 센서
  \item \textbf{화면 출력이 필요한 경우}: OLED 디스플레이 추가
\end{itemize}
\end{keypoint}

\newpage

\subsection{Grove I2C 연결 방식}

Grove 센서는 \textbf{Grove Shield}를 통해 마이크로비트에 연결됩니다.

\begin{figure}[H]
\centering
\includegraphics[width=0.7\textwidth]{connect_arduino.jpg}
\caption{Grove Shield와 센서 연결 예시}
\end{figure}

\subsubsection{I2C 통신이란?}

\textbf{I2C}(Inter-Integrated Circuit)는 \textbf{2개의 선(SDA, SCL)}으로 여러 장치를 연결하는 통신 방식입니다.

\begin{keypoint}
\textbf{I2C 통신의 특징}
\begin{itemize}
  \item \textbf{2개 선 공유}: SDA(데이터), SCL(클럭)
  \item \textbf{고유 주소}: 각 센서는 고유 주소로 구분 (예: 0x44, 0x76)
  \item \textbf{다중 연결}: 이론상 127개, 실용적으로 5~8개 장치 연결 가능
  \item \textbf{간단한 배선}: 납땜 없이 Grove 커넥터로 쉽게 연결
\end{itemize}
\end{keypoint}

\subsubsection{연결 구조}

\begin{center}
\begin{tikzpicture}[
  node distance=1.5cm,
  component/.style={font=\small}
]

% 마이크로비트
\node[rectangle, rounded corners=5pt, minimum width=3cm, minimum height=1.5cm,
      fill=microbitcolor!20, draw=microbitcolor, line width=2pt] (microbit) {
  \textbf{micro:bit V2}
};

% 연결 화살표
\draw[-Stealth, line width=2pt, mainblue] (microbit.south) -- ++(0,-0.5);

% Grove Shield
\node[rectangle, rounded corners=5pt, minimum width=5cm, minimum height=1.2cm,
      fill=mainblue!20, draw=mainblue, line width=2pt, below=2cm of microbit] (shield) {
  \textbf{Grove Shield}
};

\node[below=0.2cm of shield, font=\footnotesize, text=mainblue] {최대 6개 I2C 포트};

% I2C 포트들
\foreach \x in {-1.8, -1, -0.2, 0.6, 1.4, 2.2} {
  \node[rectangle, rounded corners=2pt, minimum width=0.5cm, minimum height=0.3cm,
        fill=yellow!70, draw=black] at ([yshift=-0.5cm, xshift=\x cm]shield.south) {};
}

% 센서들
\foreach \x/\name/\addr in {-1.8/SHT40/0x44, -0.2/BMP280/0x76, 1.4/SGP30/0x58} {
  \draw[-Stealth, line width=1.5pt, red] ([yshift=-0.5cm, xshift=\x cm]shield.south) -- ++(0,-0.5);
  \node[rectangle, rounded corners=3pt, minimum width=1cm, minimum height=0.6cm,
        fill=red!20, draw=red, line width=1pt, 
        below=1.2cm of shield, xshift=\x cm, font=\tiny] {
    \textbf{\name}\\[-0.1em]\textcolor{gray}{\addr}
  };
}

% I2C 선 표시
\draw[dashed, mainblue, line width=1pt] 
  ([xshift=-2.3cm]shield.west) -- ([xshift=2.7cm]shield.east);
\node[left, font=\tiny, text=mainblue] at ([xshift=-2.3cm]shield.west) {SDA/SCL};

\end{tikzpicture}
\end{center}

\subsection{센서 조합 가능성}

마이크로비트 내장 센서 6개 + 데이터로거 1개 + Grove 확장 센서 최대 6개를 조합하면 \textbf{총 848가지 조합}이 가능합니다!

\begin{table}[H]
\centering
\begin{tabularx}{\textwidth}{|c|c|X|}
\hline
\rowcolor{lightgray}
\textbf{외장 센서 수} & \textbf{조합 수} & \textbf{데이터 수 범위} \\
\hline
0개 & 1 & 10개 (내장만) \\
\hline
1개 & 10 & 11~14개 \\
\hline
2개 & 45 & 12~18개 \\
\hline
3개 & 120 & 13~22개 \\
\hline
4개 & 210 & 14~25개 \\
\hline
5개 & 252 & 15~27개 \\
\hline
6개 & 210 & 16~28개 \\
\hline
\textbf{총합} & \textbf{848} & \textbf{10~28개} \\
\hline
\end{tabularx}
\caption{센서 조합 통계}
\end{table}

\newpage

\subsection{센서 활용 예시}

\subsubsection{학교 현장 적용 프로젝트}

\begin{table}[H]
\centering
\begin{tabularx}{\textwidth}{|l|l|X|}
\hline
\rowcolor{lightgray}
\textbf{프로젝트 주제} & \textbf{사용 센서} & \textbf{측정 목적} \\
\hline
교실 환경과 집중도 & 온도, 소음, 빛 & 쾌적한 학습 환경 조건 찾기 \\
\hline
복도 소음 분석 & 소음 & 시간대별 소음 패턴 파악 \\
\hline
급식실 대기 시간 & 거리 & 줄 서는 시간 측정 \\
\hline
운동장 활동량 & 가속도 & 체육 시간 활동량 측정 \\
\hline
도서관 이용 패턴 & 빛, 소음 & 조용한 시간대 분석 \\
\hline
교실 공기질 관리 & CO₂, 온습도 & 환기 시점 파악 \\
\hline
\end{tabularx}
\caption{센서 활용 프로젝트 예시}
\end{table}

\begin{warning}
\textbf{센서의 한계 이해하기}

센서는 \textbf{현재 상태}만 측정할 뿐입니다:
\begin{itemize}
  \item ❌ 센서는 \textbf{미래를 예측하지 않습니다}
  \item ❌ 센서는 \textbf{원인을 알려주지 않습니다}
  \item ❌ 센서는 \textbf{판단하지 않습니다}
  \item ✅ 센서는 \textbf{지금 이 순간의 값}만 알려줍니다
\end{itemize}

예: 온도 센서가 "28도"를 측정했다면
\begin{itemize}
  \item ✅ "지금 온도는 28도다" (사실)
  \item ❌ "더워서 학생들이 집중을 못한다" (추정)
\end{itemize}
\end{warning}

\subsection{센서 탐색 활동}

\subsubsection{활동 1: 센서 시뮬레이터 체험}

프로젝트에 포함된 \textbf{센서 시뮬레이터}를 사용하여:
\begin{itemize}
  \item 드래그 앤 드롭으로 센서 연결 체험
  \item 각 센서가 측정하는 데이터 확인
  \item 센서 조합 가능성 탐색
\end{itemize}

\subsubsection{활동 2: 센서 검색 도구 활용}

\textbf{AI 센서 추천 도구}를 사용하여:
\begin{itemize}
  \item "교실 공기가 탁한지 알고 싶어요" 입력
  \item AI가 필요한 센서 추천
  \item 추천 이유 확인
\end{itemize}

\subsubsection{활동 3: 실물 센서 관찰}

교수자가 준비한 센서들을 직접 보고 만져보며:
\begin{itemize}
  \item 센서의 크기와 형태 확인
  \item Grove 커넥터 연결 방식 이해
  \item 센서별 특징 파악
\end{itemize}

\newpage

%===========================================
\section{프로젝트 구상}
%===========================================

\subsection{프로젝트 주제 브레인스토밍}

3주차부터는 본격적으로 개인 프로젝트를 시작합니다. 이번 주에는 \textbf{프로젝트 아이디어를 구상}해봅니다.

\subsubsection{생각해볼 질문들}

\begin{itemize}
  \item 학교에서 \textbf{무엇이 궁금}한가요?
  \item 어떤 \textbf{문제를 해결}하고 싶나요?
  \item 어떤 \textbf{환경 요소}가 영향을 줄까요?
  \item \textbf{어디서 데이터}를 모을 수 있나요? (실습 학교)
\end{itemize}

\subsubsection{좋은 주제의 조건}

\begin{keypoint}
\textbf{좋은 프로젝트 주제란?}
\begin{itemize}
  \item \textbf{측정 가능}: 센서로 측정할 수 있는 것
  \item \textbf{실습 가능}: 실습 학교에서 데이터 수집 가능
  \item \textbf{의미 있음}: 초등교육에 도움이 되는 것
  \item \textbf{실현 가능}: 4주 안에 완성할 수 있는 것
\end{itemize}
\end{keypoint}

\subsection{피해야 할 주제}

\begin{warning}
\textbf{이런 주제는 피하세요}
\begin{itemize}
  \item 너무 추상적: "학생들의 행복도 측정"
  \item 측정 불가능: "학생들의 창의성 향상"
  \item 데이터 수집 곤란: "학생 집의 환경 조사"
  \item 시간 부족: "1년 동안의 계절별 변화 분석"
\end{itemize}
\end{warning}

\subsection{다음 주 준비}

3주차에는 다음을 준비해야 합니다:
\begin{enumerate}
  \item 프로젝트 주제 최종 결정
  \item 필요한 센서 선택
  \item 데이터 수집 계획서 작성
\end{enumerate}

미리 생각해두면 3주차 수업이 더 수월합니다!

\newpage

%===========================================
\section{과제 안내}
%===========================================

\subsection{이번 주 과제}

\begin{tcolorbox}[enhanced, colback=blue!5, colframe=NavyBlue, 
                  title={\Large\textbf{2주차 과제}}, fonttitle=\bfseries]

\textbf{과제 1: 바이브코딩 체험}
\begin{itemize}
  \item Claude를 사용하여 자기소개 웹페이지 만들기
  \item 필수 포함: 이름, 전공, 관심사
\end{itemize}

\vspace{0.5em}

\textbf{과제 2: 센서 탐색}
\begin{itemize}
  \item 센서 시뮬레이터로 다양한 센서 탐색
  \item 관심 있는 센서 3개 이상 조사
\end{itemize}

\end{tcolorbox}

\subsection{제출물}

\begin{table}[H]
\centering
\begin{tabularx}{\textwidth}{|l|X|}
\hline
\rowcolor{lightgray}
\textbf{제출 항목} & \textbf{내용} \\
\hline
1. 결과물 (HTML) & 완성된 자기소개 웹페이지 파일 \\
\hline
2. 대화목록 전문 (PDF) & Claude와 주고받은 전체 대화 내용 \\
\hline
3. 대화목록 요약본 (PDF) & 주요 프롬프트와 수정 과정 요약 \\
\hline
\end{tabularx}
\end{table}

\subsubsection{대화목록 PDF 생성 방법}

Claude 대화 페이지에서:
\begin{enumerate}
  \item 대화 내용 전체 선택 (Ctrl+A)
  \item 복사 (Ctrl+C)
  \item 워드나 구글 문서에 붙여넣기
  \item PDF로 저장
\end{enumerate}

또는 Claude에게 직접 요청:
\begin{verbatim}
이 대화를 PDF로 만들 수 있도록 정리해줘
\end{verbatim}

\subsection{제출 방법 및 기한}

\begin{itemize}
  \item \textbf{제출 기한}: 2월 2일 (일) 자정까지
  \item \textbf{제출 방법}: LMS 과제 게시판에 업로드
  \item \textbf{파일명 형식}:
  \begin{itemize}
    \item 결과물: \texttt{2주차\_이름\_결과물.html}
    \item 대화목록 전문: \texttt{2주차\_이름\_대화전문.pdf}
    \item 대화목록 요약: \texttt{2주차\_이름\_대화요약.pdf}
  \end{itemize}
\end{itemize}

\subsection{평가 기준}

\begin{table}[H]
\centering
\begin{tabularx}{\textwidth}{|l|c|X|}
\hline
\rowcolor{lightgray}
\textbf{평가 항목} & \textbf{배점} & \textbf{평가 내용} \\
\hline
결과물 완성도 & 40\% & 필수 내용 포함, 동작 확인 \\
\hline
프롬프트 작성 & 30\% & 구체성, 단계적 수정 과정 \\
\hline
대화목록 정리 & 20\% & 주요 내용 요약, 가독성 \\
\hline
창의성 & 10\% & 독창적인 아이디어, 디자인 \\
\hline
\end{tabularx}
\end{table}

\begin{keypoint}
\textbf{과제 수행 팁}
\begin{itemize}
  \item 조금씩 단계적으로 수정하기
  \item 각 단계마다 결과 확인하기
  \item 막히면 교수님이나 친구에게 질문하기
  \item 대화 기록을 꼭 저장하기!
\end{itemize}
\end{keypoint}

\newpage

%===========================================
\section{참고 자료}
%===========================================

\subsection{유용한 링크}

\begin{itemize}
  \item \textbf{Claude}: \url{https://claude.ai}
  \item \textbf{MakeCode}: \url{https://makecode.microbit.org}
  \item \textbf{마이크로비트}: \url{https://microbit.org/ko/}
  \item \textbf{Grove 위키}: \url{https://wiki.seeedstudio.com}
  \item \textbf{강의 자료실}: LMS 참고
\end{itemize}

\subsection{추가 학습 자료}

\begin{table}[H]
\centering
\begin{tabularx}{\textwidth}{|l|X|}
\hline
\rowcolor{lightgray}
\textbf{주제} & \textbf{추천 자료} \\
\hline
바이브코딩 & Anthropic 공식 문서, Claude 사용 가이드 \\
\hline
마이크로비트 & 마이크로비트 한글 가이드, MakeCode 튜토리얼 \\
\hline
Grove 센서 & Seeed Studio Wiki, 센서 데이터시트 \\
\hline
데이터 분석 & 엑셀 기초, 그래프 그리기 \\
\hline
\end{tabularx}
\end{table}

\subsection{FAQ (자주 묻는 질문)}

\textbf{Q1: Claude가 이상한 코드를 만들어요}
\begin{itemize}
  \item A: 더 구체적으로 설명하거나, 잘못된 부분을 지적하여 수정 요청하세요
\end{itemize}

\textbf{Q2: HTML 파일이 실행되지 않아요}
\begin{itemize}
  \item A: 파일 확장자가 .html인지 확인하고, 브라우저로 열어보세요
\end{itemize}

\textbf{Q3: 센서를 어떻게 선택해야 하나요?}
\begin{itemize}
  \item A: 프로젝트 주제에 맞는 센서를 선택하세요. 3주차에 자세히 다룹니다
\end{itemize}

\textbf{Q4: 대화목록이 너무 길어요}
\begin{itemize}
  \item A: 요약본에는 주요 프롬프트와 결과만 포함하면 됩니다
\end{itemize}

\newpage

%===========================================
\section{정리 및 다음 주 예고}
%===========================================

\subsection{핵심 요약}

\begin{keypoint}
\textbf{2주차 핵심 내용}

\begin{enumerate}
  \item \textbf{바이브코딩 = AI와 대화하며 코드 만들기}
  \begin{itemize}
    \item 문제 정의가 가장 중요
    \item 효과적인 프롬프트 작성법 익히기
  \end{itemize}
  
  \item \textbf{Claude 사용법}
  \begin{itemize}
    \item 구체적으로 요청하기
    \item 단계적으로 수정하기
  \end{itemize}
  
  \item \textbf{센서 탐색}
  \begin{itemize}
    \item 내장 6개 + Grove 확장 센서
    \item 848가지 조합 가능
    \item 센서는 현재 상태만 측정
  \end{itemize}
  
  \item \textbf{프로젝트 구상 시작}
  \begin{itemize}
    \item 학교 현장 적용 가능한 주제
    \item 측정 가능하고 의미 있는 것
  \end{itemize}
\end{enumerate}
\end{keypoint}

\subsection{3주차 예고}

\begin{tcolorbox}[enhanced, colback=green!5, colframe=ForestGreen, 
                  title={\Large\textbf{3주차: 프로젝트 기획}}, fonttitle=\bfseries]

\textbf{주요 내용}
\begin{itemize}
  \item 프로젝트 주제 확정
  \item 필요한 센서 선택
  \item 데이터 수집 계획서 작성
  \item 프로젝트 예시 분석
\end{itemize}

\vspace{0.5em}

\textbf{준비사항}
\begin{itemize}
  \item 관심 있는 프로젝트 주제 3개 정도 생각해오기
  \item 어떤 센서가 필요할지 고민해보기
  \item 실습 학교에서 어디서 데이터를 수집할지 생각해보기
\end{itemize}

\end{tcolorbox}

\vfill

\begin{center}
\Large\textbf{수고하셨습니다!}

\vspace{1em}

\normalsize
다음 주에 만나요 😊
\end{center}

\newpage

%===========================================
% 문서 정보
%===========================================

\section*{문서 정보}
\addcontentsline{toc}{section}{문서 정보}

\begin{table}[H]
\centering
\begin{tabularx}{0.8\textwidth}{|l|X|}
\hline
\rowcolor{lightgray}
\textbf{항목} & \textbf{내용} \\
\hline
문서명 & 사물인터넷과 빅데이터 2주차 강의자료 \\
\hline
버전 & v1.1.0.0 \\
\hline
작성일 & 2026-01-28 \\
\hline
작성자 & 양창모 \\
\hline
소속 & 청주교육대학교 컴퓨터교육과 \\
\hline
페이지 수 & \pageref{LastPage} 페이지 \\
\hline
\end{tabularx}
\end{table}

\subsection*{생성/수정 이력}

\begin{table}[H]
\centering
\begin{tabularx}{\textwidth}{|c|c|c|c|X|}
\hline
\rowcolor{lightgray}
\textbf{버전} & \textbf{날짜} & \textbf{시간} & \textbf{수준} & \textbf{변경 내용} \\
\hline
v1.0.0.0 & 2026-01-28 & 14:00 & 최초 & 초안 작성 (교재 형식, 30쪽) \\
\hline
v1.1.0.0 & 2026-01-28 & 15:30 & 마이너 & Grove 센서 10개 이미지 및 상세 설명 추가 \\
\hline
\end{tabularx}
\end{table}

\vfill

\begin{center}
\rule{0.5\textwidth}{0.4pt}\\[0.5em]
\textit{청주교육대학교 컴퓨터교육과}\\
\textit{사물인터넷과 빅데이터 --- 2주차 강의자료}\\
\textit{v1.1.0.0 | 2026-01-28}
\end{center}

\end{document}